\section{Introduction}

\subsection{Background and Significance}

Joint-Embedding Predictive Architectures (JEPAs) have become a leading recipe for
model-based control from high-dimensional observations: an encoder maps observations to a
compact latent space, a predictor models the latent dynamics, and a planner---most often
model predictive control (MPC)---optimises a short action sequence by rolling the predictor
forward and minimising a latent goal distance \cite{lecun-2022-path,zhou-2025-dinowm,
sobal-2025-offline,wang-2026-temporal,assran-2025-vjepa2}. The appeal is that planning never
touches pixels, so the model is cheap, and that training needs no reward labels.

In practice the world model is \emph{frozen} after training, and its usefulness at deployment
is decided entirely by how well the training distribution covers the test environment. Yet
test environments shift: objects change shape, cameras change colour balance, noise appears,
masses and frictions change, and room layouts differ. Under such shifts latent rollouts
drift, MPC optimises the wrong objective, and success collapses---even for high-capacity
models \cite{zhou-2025-dinowm,toso-2026-invariant,zhang-2026-hierarchical}.

AdaJEPA \cite{wang-2026-adajepa} made a strong and useful observation: the world model does
not have to stay frozen. Inside the MPC loop, the transition that results from every executed
action is a free self-supervised label. AdaJEPA plans, executes an action chunk, takes one
gradient step on the predictor's last transformer block and the encoder head using the
observed transition, and replans with the updated model. One step per replan already recovers
a large fraction of the loss under visual and dynamics shifts, at negligible latency.

\subsection{The limitation we address}

AdaJEPA's adaptation is \emph{uniform}: the same one gradient step, at the same learning
rate, applied at the same layers, whether the model already predicts the current environment
well or is being asked about dynamics it has never seen. Three consequences follow.

\begin{enumerate}[leftmargin=1.2em,itemsep=2pt]
\item \textbf{Misallocated compute.} In-distribution episodes are updated as hard as
out-of-distribution ones. Adaptation is not free: extra gradient steps can \emph{hurt} when
the model is already adequate, because they fit the model to a five-transition buffer and move
it off the pretrained manifold.
\item \textbf{Misallocated plasticity.} Directly updating the pretrained predictor weights
offers no structural guarantee that the update is local; the setting where adaptation matters
most (unseen dynamics) is also the setting where a large step is most damaging.
\item \textbf{No memory across episodes.} AdaJEPA's evaluation---and its design---reset the
model at every episode. Everything learned about an environment is discarded when the episode
ends, so nothing accumulates: the model can never get better at a \emph{family} of
environments, only at the current one.
\end{enumerate}

A fourth, subtler issue is architectural. A single feed-forward predictor spends the same
amount of computation on every prediction. Recent work on reasoning with small networks
suggests the opposite is better: improve an answer by recursing on it in latent space, with
deep supervision at every recursion depth, and let the number of recursions depend on
difficulty \cite{tiny-recursive-2025,hrm-2025,hrm-text-2026}. Difficulty-dependent compute is
exactly the missing knob for a world model under shift.

\subsection{Approach and contributions}

We propose \textbf{DTA-JEPA}, a dual-timescale recursive adaptive latent world model. Our
contributions are:

\begin{enumerate}[leftmargin=1.2em,itemsep=2pt]
\item \textbf{Computational dual timescale (\S\ref{sec:model}).} The predictor is split into a
slow \emph{S-module} that turns a short latent--action history into a dynamics context $c_t$,
and a fast \emph{R-module} that recursively refines the next latent
$\yhat^{(k)} = \yhat^{(k-1)} + g_{\f}([\yhat^{(k-1)},c_t],u_t)$ for $k=1,\dots,K_t$, with deep
supervision over all depths. The slow module is the pretrained anchor and is never modified
online; all plasticity lives in a rank-4 LoRA fast parameter set attached to the R-module and
the encoder head.
\item \textbf{Uncertainty as a resource allocator (\S\ref{sec:unc}).} A single scalar---the
Mahalanobis surprise of the observed transition under the model's own predicted
Gaussian---drives the number of gradient steps $U_t$, the learning-rate multiplier, the
refinement depth $K_t$, the planning penalty, the memory write priority and the rollback
decision. Compute is therefore spent where the model is measurably wrong, and saved where it
is right.
\item \textbf{Dual memory and safe consolidation (\S\ref{sec:mem}).} An episodic memory of
high-surprise latent transitions (persistent across episodes) plus a fingerprint-indexed bank
of fast-parameter snapshots give cross-episode accumulation without catastrophic forgetting;
an anchor ball on $\f$, a validation-gated rollback and a latent support penalty in the
planning cost bound the damage any single update can do. Slow parameters are consolidated
only at episode boundaries under an EWC penalty.
\item \textbf{Protocol-level replication and empirical comparison (\S\ref{sec:setup},
\S\ref{sec:results}).} We replicate AdaJEPA's four shift families (shape, visual, dynamics,
layout) plus a continual environment-switching suite on a self-contained benchmark, and
compare Frozen, a faithful AdaJEPA re-implementation and DTA-JEPA under an identical planner,
horizon, replay buffer and episode set, together with five ablations.
\end{enumerate}

\subsection{What we do not claim}

Our benchmark is not PushT or MuJoCo. It is a physics-based 2D re-implementation built to make
the \emph{evaluation protocol} comparable---the same shift families, the same segment-based
goal sampling, the same MPC budget, the same success criteria---rather than the same pixels.
We therefore report relative comparisons between methods evaluated under identical conditions,
not absolute numbers transferable to the original benchmark.
